
\newtcolorbox{notizTemplate}[1][]{
	enhanced jigsaw,
	lines before break=9,
	breakable,
	parbox=false,
	attach boxed title to top text left={yshift=-3mm,yshifttext=-3mm},
	title=\emoji{pen} Notizen,
	colback=white,
	colframe=ForestGreen,
	colbacktitle=ForestGreen,
	#1}

\newcommand{\notiz}[1]{\begin{notizTemplate}[height=#1cm]
\end{notizTemplate}}
	
\section{Einführung}
%\begin{handwritenotes}[Write your thoughts or examples below:]
\subsection{Definition der Stadt}
\begin{enumerate}
\item \textbf{Statistische Definition}
\begin{figure}[H]
\centering
\includegraphics[width=0.8\linewidth]{"Private/Images/surseePopulation"}
\caption{Bevölkerungswachstum der Stadt Sursee über die Zeit}
\end{figure}
\begin{multicols}{2}
\notiz{4}
\columnbreak
\begin{table}[H]
\centering
\begin{tabular}{c|c}
\textbf{Stadtdefinition} & \textbf{Schwelle} \\\midrule\midrule
Island, Dänemark & 200 \\\midrule
Deutschland, Frankreich & 2'000 \\\midrule
Schweiz & 10'000 \\\midrule
Japan & 50'000 \\
\end{tabular}
\end{table}
\end{multicols}
\item \textbf{Historische Definition}
\begin{figure}[H]
\centering
\includegraphics[width=0.5\linewidth]{"Private/Images/werdenbergSG"}
\caption{Stadt Werdenberg SG}
\end{figure}
\begin{figure}[H]
\centering
\includegraphics[width=\linewidth]{"Private/Images/histBegriffTemplate"}
\caption{Dimensionen historischer Städte nach \textcite{hep}}
\end{figure}
\begin{multicols}{2}
\notiz{4}
\columnbreak
\begin{table}[H]
\centering
\textbf{Stadt-Rechte}
\begin{enumerate}
\item Münzrecht
\item Marktrecht
\item Befestigungsrecht
\item etc.
\end{enumerate}
\end{table}
\end{multicols}
\item \textbf{Funktionale Definition}
\begin{figure}[H]
\centering
\includegraphics[width=0.8\linewidth]{"Private/Images/funktionsraum"}
\caption{Funktionale Aspekte einer Stadt (Paris: Aussicht vom \textit{Arc de Triomphe})}
\end{figure}
\begin{multicols}{2}
\notiz{4}
\columnbreak
\begin{table}[H]
\centering
\textbf{7 Grunddaseinsfunktionen}
\begin{enumerate}
\item Arbeiten
\item Wohnen
\item Sich erholen
\item Sich bilden
\item (Sich) versorgen (und entsorgen)
\item An gemeinschaftlichen Leben teilnehmen
\item Sich fortbewegen
\end{enumerate}
\end{table}
\end{multicols}
\item Geografische Definition \begin{multicols}{2}
\notiz{4}
\columnbreak

\begin{figure}[H]
\centering
\includegraphics[width=\linewidth]{"Private/Images/geoStadtbegriff"}
\caption{Geographischer Stadtbegriff nach \textcite{westermann}}
\end{figure}

\end{multicols}
\item Subjektive Definition \notiz{4}
\end{enumerate}
\newpage
\textbf{Begriffe}
\begin{itemize}
\item Agglomeration 
\begin{figure}[H]
\centering
\includegraphics[width=0.7\linewidth]{"Private/Images/agglomeration"}
\caption{Agglomerationsgebiete der Schweiz (Ausschnitt), \href{https://www.atlas.bfs.admin.ch/maps/13/map/mapIdOnly/27623_de.html}{Quelle}}
\end{figure}
\notiz{4}
\item Raum mit städtischem Charakter \notiz{4}
\item Megastadt \notiz{4}
\item Metropole \notiz{4}
\item Metropol-Region 
\begin{multicols}{2}
\notiz{4}
\columnbreak
\begin{figure}[H]
\centering
\includegraphics[width=\columnwidth]{"Private/Images/gza"}
\caption{Mitgliedskantone der ``Greater Zurich Area''}
\end{figure}
\end{multicols}
\item Global City \notiz{4}
\item Megalopolis 
\begin{multicols}{2}
\notiz{4}
\columnbreak
\begin{figure}[H]
\centering
\includegraphics[width=\columnwidth]{"Private/Images/megalopolis"}
\caption{Megalopolis in den USA}
\end{figure}
\end{multicols}
\end{itemize}


\begin{myexercise}
Suchen Sie sich in 4er-Gruppen eine Stadt aus, die Sie im Rahmen der nächsten 7 Doppellektionen vertiefen werden:
\begin{itemize}
\item Basel (\& Laufen)
\item Bern (\& Huttwil)
\item Genf (\& Allaman)
\item Andere Stadt (\& Dorf) nach Wunsch
\end{itemize}
Erarbeiten Sie eine Kurzpräsentation auf nächstes Mal, in welcher Sie die statistische, historische, formale, geografische, subjektive und funktionalen Definitionen dieser Städte ausarbeiten (die ``Dörfer'' in den Klammern werden wir erst später anschauen).

Bereiten Sie auf nächstes Mal eine Kurzpräsentationen à 5 Minuten vor, in der Sie folgendes darstellen:
\begin{itemize}
\item \textbf{Statistische Definition}: Einwohnerzahl
\item \textbf{Historische Definition}: Historische Entstehung, Spezialrechte (kurz)?
\item \textbf{Geografische Definition}: Welche Aspekte sind besonders offensichtlich? Welche weniger?
\item \textbf{Funktionale Definition}: Suchen Sie ein Luftbild Ihrer Stadt und annotieren Sie darauf die 7 Grunddaseinsfunktionen!
\end{itemize}

Link Luftbilder: \qrcode{https://www.lds.geo.uzh.ch/de/Alle-Orte.html}

Formales: 5‘ Präsentation mit mindestens 2 Bildern. Form frei (PPT / Poster / Sonstiges). Verwendung von ChatGPT erlaubt, Ihre Präsentation sollte aber korrekt sein!


\end{myexercise}
\newpage
\begin{mychallenge}[Ausschnitt aus Buch ``Klett Fundamente'' zu Stadtbegriffen]
\foreach \n in {1,...,3}{
\includegraphics[width=\linewidth, page=\n]{"Private/Appendix/Der geographische Stadtbegriff.pdf"}\\
}


\end{mychallenge}


\begin{mychallenge}[Verstädterung in der Schweiz \cite{geoaktiv}]
\includegraphics[width=\linewidth]{"Private/Appendix/verstaedterung"}


\end{mychallenge}